\scanpage{170}
\mspage{86}
% archive sheet number: 169
\[
\psi_J:I\longrightarrow\underline{\operatorname{Hom}}(J,C)
\]
% lines 2-5
comme un foncteur
\[
I\times J\xrightarrow{\psi}C,
\]
on trouve la ``représentation indicielle'' de $\varphi$
\[
\varphi(j)=\underbrace{\varinjlim_{i\in I}\underbrace{\psi(i,j)}_{\deleted{\unsure{\in\operatorname{Ob}C}}}}_{\substack{\text{limite inductive}\\\text{prise dans }\operatorname{Ind}'_{\pi}(C)\\\text{d'objets de }C\subset\operatorname{Ind}'_{\pi}(C)}},
\]

% lines 6-9
d'où
\[
\begin{aligned}
\underbrace{\varinjlim_{j\in J}{}_{\operatorname{Ind}'_{\pi}C}\varphi(j)}_{\substack{\text{sous réserve}\\\text{d'existence}}}
&=\varinjlim_{j\in J}\varinjlim_{i\in I}\psi(i,j)\\
&=\varinjlim_{i\in I}\varinjlim_{j\in J}\psi(i,j),
\end{aligned}
\]
% lines 10-21
toutes les limites inductives étant prises dans $\operatorname{Ind}'_{\pi}(C)$. Mais \deleted{si \unsure{pour}} \inserted{les} $\varinjlim$ de type $J$ sont représentables dans $C$, et $\operatorname{card}J\leq\pi$, on sait que \inserted{la} $\varinjlim_{j\in J}\psi(i,j)$, \uline{dans $C$} est aussi leur $\varinjlim$ dans $\operatorname{Ind}'_{\pi}C$, et on trouve donc que $\varinjlim_{j\in J}\varphi(j)$ existe dans $\operatorname{Ind}'_{\pi}(C)$ sss $\varinjlim_{i\in I}\Psi(i)$ existe dans $\operatorname{Ind}'_{\pi}C$, où pour $i\in I$, on pose
\[
\Psi(i)=\varinjlim_{j\in J}\psi(i,j)\quad\text{\uline{dans $C$}}.
\]
